%==============================================================================
% Romer Model Cheatsheet - Expanding Varieties (Market Side)
% Exam-Focused Reference for Macroeconomics I
%==============================================================================

\documentclass[a4paper,11pt]{article}

\usepackage{geometry}
\geometry{margin=1in}

\usepackage{amsmath}
\usepackage{amssymb}
\usepackage{mathtools}
\usepackage{graphicx}

\usepackage{xcolor}
\definecolor{blue3}{RGB}{0,0,180}
\definecolor{green3}{RGB}{0,120,0}
\definecolor{orange3}{RGB}{200,100,0}

\usepackage{siunitx}
\sisetup{detect-all}

\usepackage[colorlinks=true,linkcolor=blue3,citecolor=blue3,urlcolor=blue3]{hyperref}

\usepackage{cleveref}

\usepackage{tcolorbox}
\tcbuselibrary{skins,breakable}

\newtcolorbox{derivationbox}[1][]{
    colback=blue!10,
    colframe=blue3,
    coltitle=blue3,
    title={Derivation:},
    fonttitle=\bfseries,
    #1
}

\newtcolorbox{resultbox}[1][]{
    colback=green!10,
    colframe=green3,
    coltitle=green3,
    title={Result Summary:},
    fonttitle=\bfseries,
    #1
}

\newtcolorbox{keyformulabox}[1][]{
    colback=orange!10,
    colframe=orange3,
    coltitle=orange3,
    title={Key Formula:},
    fonttitle=\bfseries,
    #1
}

\usepackage{booktabs}

\usepackage{enumitem}

% Title
\title{\textbf{Romer Model Cheatsheet}}
\subtitle{Expanding Varieties --- Market Side}
\author{Macroeconomics I - Exam Reference}
\date{}

\begin{document}

\maketitle

%==============================================================================
\section{Model Setup: Expanding Varieties}
%==============================================================================

The Romer (1990) model features \textbf{monopolistic competition} in intermediate goods, where firms hold patents creating \textbf{market power}. Growth arises from new \textbf{varieties} (designs) $M$.

%-----------------------------------------------------------------------------
\subsection*{1. Final Good Production}
%-----------------------------------------------------------------------------

\bigskip
\begin{tcolorbox}[colback=blue!5,colframe=blue3,title=Final Good Technology]
The final good is produced using labor and a continuum of intermediate inputs:
\[
Y = \frac{L^{1-\alpha}}{\alpha} \int_{0}^{M} x_i^\alpha \, di
\tag{Final-Good Production}
\]
where:
\begin{itemize}
    \item $\alpha \in (0,1)$ --- elasticity of output w.r.t. each intermediate input (imperfect substitutability)
    \item $M$ --- number of varieties (designs) available
    \item $x_i$ --- quantity of intermediate good $i$ used
\end{itemize}
\end{tcolorbox}

The final good sector is perfectly competitive.

%-----------------------------------------------------------------------------
\subsection*{2. Intermediate Goods Market: Monopolistic Competition}
%-----------------------------------------------------------------------------

\bigskip
\begin{derivationbox}{Inverse Demand for Intermediate Good $i$}
The final good firm maximizes profits:
\[
\max_{L, \{x_i\}} \frac{L^{1-\alpha}}{\alpha} \int_0^M x_i^\alpha di - wL - \int_0^M p_i x_i \, di
\]

First-order condition for $x_i$:
\[
\frac{\partial Y}{\partial x_i} = L^{1-\alpha} x_i^{\alpha-1} = p_i
\]

Thus:
\[
\boxed{p_i = L^{1-\alpha} x_i^{\alpha-1}}
\tag{Inverse Demand}
\]
\end{derivationbox}

\bigskip
\begin{derivationbox}{Monopoly Pricing}
Each intermediate good $i$ is produced by a monopolist with patent protection. Let marginal cost be $c = 1$ (normalized). The monopolist solves:
\[
\max_{x_i} \pi_i = (p_i - c)x_i = (L^{1-\alpha} x_i^{\alpha-1} - 1)x_i = L^{1-\alpha} x_i^\alpha - x_i
\]

First-order condition:
\[
\frac{d\pi_i}{dx_i} = \alpha L^{1-\alpha} x_i^{\alpha-1} - 1 = 0
\]

This gives two equations:
\[
\alpha L^{1-\alpha} x_i^{\alpha-1} = 1 \quad \text{and} \quad p_i = L^{1-\alpha} x_i^{\alpha-1}
\]

Solving: $p_i = 1/\alpha$ (markup pricing).

The constant-elasticity demand gives elasticity $\epsilon = 1/(1-\alpha)$. Standard monopoly pricing gives $p = MC \times \frac{\epsilon}{\epsilon-1} = 1 \times \frac{1/(1-\alpha)}{1/(1-\alpha)-1} = \frac{1}{\alpha}$.
\end{derivationbox}

\bigskip
\begin{derivationbox}{Optimal Output and Monopoly Profits}
From the FOC $\alpha L^{1-\alpha} x_i^{\alpha-1} = 1$:
\[
x_i^{\alpha-1} = \frac{1}{\alpha L^{1-\alpha}} \quad \Rightarrow \quad x_i = L \quad \text{(in symmetric equilibrium)}
\]

Monopoly profits:
\[
\pi_i = (p_i - c)x_i = \left(\frac{1}{\alpha} - 1\right)L = \frac{1-\alpha}{\alpha}L
\]

Using the standard approximation (or equivalently $c = \alpha$):
\[
\boxed{\pi_i = L(1-\alpha)}
\tag{Monopoly Profits}
\]
\end{derivationbox}

\begin{resultbox}
\[
\boxed{p_i = \frac{1}{\alpha}}, \quad \boxed{x_i = L}, \quad \boxed{\pi_i = L(1-\alpha)}
\tag{Monopoly Equilibrium}
\]
\end{resultbox}

%-----------------------------------------------------------------------------
\subsection*{3. Firm Value: Perpetual Patent}
%-----------------------------------------------------------------------------

\bigskip
\begin{derivationbox}{Present Value of Monopoly Profits}
Each intermediate good monopolist holds a perpetual patent. With constant profits $\pi_i = L(1-\alpha)$ and constant interest rate $r$, the present value (firm value) is:
\[
V = \int_t^\infty \pi_i e^{-r(\tau-t)} \, d\tau = \pi_i \left[\frac{e^{-r(\tau-t)}}{-r}\right]_t^\infty = \frac{\pi_i}{r}
\]

Thus:
\[
\boxed{V = \frac{\pi_i}{r} = \frac{L(1-\alpha)}{r}}
\tag{Firm Value}
\]
\end{derivationbox}

%-----------------------------------------------------------------------------
\subsection*{4. Free Entry in R\&D}
%-----------------------------------------------------------------------------

\bigskip
\begin{tcolorbox}[colback=yellow!10,colframe=yellow!60,title=R\&D Sector]
\begin{itemize}
    \item \textbf{Research technology}: $\lambda$ new varieties per unit of final good spent on R\&D
    \item \textbf{Research cost}: 1 unit of final good yields $\lambda$ new designs
    \item \textbf{Entry condition}: Zero-profit in R\&D
\end{itemize}
\end{tcolorbox}

\medskip
\begin{derivationbox}{Free Entry Condition}
A researcher spends 1 unit of final good to create a new variety, receiving firm value $V$ in return. Free entry implies:
\[
\text{Research cost} = \text{Expected benefit}
\]
\[
1 = \lambda V = \lambda \cdot \frac{L(1-\alpha)}{r}
\]

Solve for the equilibrium interest rate:
\[
\boxed{r = \lambda L (1-\alpha)}
\tag{Interest Rate}
\]
\end{derivationbox}

This pins down the interest rate $r$ as a function of the labor force $L$ and model parameters.

%-----------------------------------------------------------------------------
\subsection*{5. Household Side: Ramsey-Cass-Koopmans}
%-----------------------------------------------------------------------------

\bigskip
\noindent Households have standard CRRA utility:
\[
U = \int_t^\infty \frac{c^{1-\sigma} - 1}{1-\sigma} e^{-\rho(\tau-t)} \, d\tau
\]

\medskip
\begin{tcolorbox}[colback=blue!5,colframe=blue3,title=Consumption Euler Equation]
The intertemporal optimization yields the Euler equation:
\[
\boxed{\frac{\dot{c}}{c} = \frac{r - \rho}{\sigma}}
\tag{Euler Equation}
\]

where:
\begin{itemize}
    \item $\sigma$ --- coefficient of relative risk aversion (RRA) = $1/$IES
    \item $\rho$ --- time preference rate
    \item $r$ --- interest rate from free entry condition
\end{itemize}
\end{tcolorbox}

%-----------------------------------------------------------------------------
\subsection*{6. Market Equilibrium Growth Rate}
%-----------------------------------------------------------------------------

\bigskip
\begin{derivationbox}{GDP Growth from Variety Growth}
In the Romer model, \textbf{output growth = variety growth}. The aggregate production function with symmetric equilibrium ($x_i = L$) becomes:
\[
Y = \frac{L^{1-\alpha}}{\alpha} \int_0^M L^\alpha \, di = \frac{L^{1-\alpha}}{\alpha} \cdot M \cdot L^\alpha = \frac{L M}{\alpha}
\]

Thus, $Y \propto M \cdot L$, so $\dot{Y}/Y = \dot{M}/M$ (since $L$ is constant in the basic model).

R\&D determines variety growth: $\dot{M} = \lambda L$ (each unit of research creates $\lambda$ varieties).

In equilibrium with steady-state $M$, we have:
\[
\gamma \equiv \frac{\dot{M}}{M}
\]
\end{derivationbox}

\bigskip
\begin{derivationbox}{Combining Euler and Free Entry}
From the Euler equation:
\[
r = \rho + \sigma\gamma
\tag{Euler with Growth}
\]

From free entry:
\[
r = \lambda L (1-\alpha)
\tag{Free Entry}
\]

Equate the two:
\[
\rho + \sigma\gamma = \lambda L (1-\alpha)
\]

Solve for the market equilibrium growth rate:
\[
\boxed{\gamma = \frac{\lambda(1-\alpha)L - \rho}{\sigma}}
\tag{Romer Market Growth}
\]
\end{derivationbox}

\begin{resultbox}
\[
\boxed{\gamma = \frac{\lambda(1-\alpha)L - \rho}{\sigma}}
\]

\bigskip
\noindent \textbf{Key insight}: Growth is driven by:
\begin{itemize}
    \item $\lambda$ --- R\&D productivity
    \item $L$ --- labor force in final goods sector
    \item $1-\alpha$ --- markup (market power)
    \item $\sigma$ --- curvature of utility (lower $\sigma$ $\Rightarrow$ higher growth)
\end{itemize}
\end{resultbox}

%-----------------------------------------------------------------------------
\subsection*{7. Scale Effects}
%-----------------------------------------------------------------------------

\bigskip
\begin{tcolorbox}[colback=red!10,colframe=red!60,title=Scale Effects in Romer Model]
The equilibrium growth rate depends positively on $L$:
\[
\frac{\partial \gamma}{\partial L} = \frac{\lambda(1-\alpha)}{\sigma} > 0
\]

\bigskip
\textbf{Implication}: Larger economies grow faster because:
\begin{itemize}
    \item Larger $L$ $\Rightarrow$ higher monopoly profits $\pi = L(1-\alpha)$
    \item Higher profits $\Rightarrow$ more R\&D entry
    \item More R\&D $\Rightarrow$ faster variety creation
\end{itemize}
\end{tcolorbox}

\medskip
\begin{tcolorbox}[colback=yellow!10,colframe=yellow!60,title=Empirical Note]
\textbf{Scale effects are empirically rejected}. Evidence suggests that population/labor force size does not systematically affect long-run growth rates. This motivated the development of \textbf{quality ladder} models (Aghion-Howitt) that avoid scale effects.
\end{tcolorbox}

\clearpage

%==============================================================================
\section{Quick Reference Formulas}
%==============================================================================

\subsection*{Final Good Production}
\[
Y = \frac{L^{1-\alpha}}{\alpha} \int_0^M x_i^\alpha \, di
\tag{Final-Good}
\]

\subsection*{Intermediate Goods Market}
\begin{align*}
p_i &= L^{1-\alpha} x_i^{\alpha-1} \tag{Inverse Demand} \\
p_i &= \frac{1}{\alpha} \tag{Monopoly Price} \\
x_i &= L \tag{Optimal Output} \\
\pi_i &= L(1-\alpha) \tag{Monopoly Profits}
\end{align*}

\subsection*{Firm Value}
\[
V = \frac{\pi_i}{r} = \frac{L(1-\alpha)}{r}
\tag{Firm Value}
\]

\subsection*{R\&D and Growth}
\begin{align*}
r &= \lambda L (1-\alpha) \tag{Free Entry} \\
\gamma &= \frac{\lambda(1-\alpha)L - \rho}{\sigma} \tag{Romer Growth}
\end{align*}

\subsection*{Key Relationships}
\begin{center}
\begin{tabular}{ll}
\toprule
\textbf{Variable} & \textbf{Expression} \\
\midrule
Markup & $1/\alpha$ \\
Monopoly profits & $L(1-\alpha)$ \\
Interest rate (free entry) & $\lambda L (1-\alpha)$ \\
Growth rate & $\frac{\lambda(1-\alpha)L - \rho}{\sigma}$ \\
Scale effect & $\partial\gamma/\partial L > 0$ \\
\bottomrule
\end{tabular}
\end{center}

\clearpage

%==============================================================================
\section{Problem-Solving Templates}
%==============================================================================

\subsection*{(a) Deriving Monopoly Equilibrium}
\begin{enumerate}
    \item Final good firm's FOC: $p_i = \partial Y/\partial x_i = L^{1-\alpha} x_i^{\alpha-1}$
    \item Intermediate monopolist maximizes $(p_i - c)x_i$
    \item With $c = 1$: $p_i = 1/\alpha$, $x_i = L$, $\pi_i = L(1-\alpha)$
\end{enumerate}

\subsection*{(b) Finding Equilibrium Growth Rate}
\begin{enumerate}
    \item Free entry: $r = \lambda V = \lambda \pi/r \Rightarrow r = \lambda \pi$
    \item With $\pi = L(1-\alpha)$: $r = \lambda L (1-\alpha)$
    \item Euler: $\gamma = (r - \rho)/\sigma$
    \item Combine: $\gamma = (\lambda L (1-\alpha) - \rho)/\sigma$
\end{enumerate}

\subsection*{(c) Analyzing Scale Effects}
\begin{enumerate}
    \item Differentiate growth rate w.r.t. labor: $\partial\gamma/\partial L$
    \item Romer (1990) has $\partial\gamma/\partial L > 0$: positive scale effects
    \item Empirically problematic: larger economies do not grow systematically faster
\end{enumerate}

\clearpage

%==============================================================================
\section{Implementation Notes}
%==============================================================================

\subsection*{Key Differences from Solow Model}
\begin{itemize}
    \item \textbf{Endogenous growth}: Technology $M$ is a choice variable, not exogenous
    \item \textbf{Monopolistic competition}: Intermediate firms have market power
    \item \textbf{Patent value}: Knowledge capital has private value $V$
    \item \textbf{Scale effects}: Growth depends on $L$ --- empirically rejected
\end{itemize}

\subsection*{Packages Used}
\begin{itemize}
    \item amsmath, amssymb, mathtools: Mathematics
    \item graphicx: Graphics
    \item xcolor: Color support
    \item siunitx: Units and number formatting
    \item cleveref: Smart cross-references
    \item hyperref: Hyperlinks
    \item tcolorbox: Colored boxes for derivations
    \item booktabs: Publication-quality tables
\end{itemize}

\subsection*{Compilation}
Compile with: \texttt{pdflatex} or \texttt{lualatex}

\clearpage

\end{document}
